% =============================================================================
% CAPÍTULO 9 — PROTOCOLOS DE DESINSTALACIÓN, LIMPIEZA Y REVERSIÓN (ROLLBACK)
% =============================================================================
\chapter[Desinstalación y Reversión]{Protocolos de Desinstalación, Limpieza y Reversión}%
\label{ch:desinstalacion}
\resumenCapitulo{Toda instalación debe poder revertirse de forma limpia y segura. Este
capítulo describe el desmantelamiento ordenado del entorno —contenedores, redes y
volúmenes de Docker— y la limpieza del estado local del cliente almacenado en el navegador.
Cierra el ciclo de vida del despliegue documentado en este manual.}

La reversión persigue dos objetivos: liberar los recursos del servidor y eliminar todo
rastro de sesión en los equipos de evaluación. El orden recomendado es: primero detener y
purgar la infraestructura de contenedores; después, limpiar el almacenamiento del cliente.
La tabla~\ref{tab:reversion} resume el alcance de cada nivel de limpieza.\par

\vspace{0.2cm}
\noindent
\renewcommand{\arraystretch}{1.35}
\fontsize{8.7}{10.4}\selectfont\sffamily
\begin{tabularx}{\dimexpr\textwidth-2pt\relax}{|>{\bfseries\color{colorPrimario}}l|X|c|}
\hline
\rowcolor{headerTabla}
\textbf{Nivel} & \textbf{Acción} & \textbf{Reversible} \\
\hline
Detención & Para los contenedores sin borrar nada. & Sí \\
\hline
\rowcolor{grisTecnico}
Purga & Elimina contenedores, redes y volúmenes. & No \\
\hline
Limpieza cliente & Borra tokens y datos del navegador. & No \\
\hline
\end{tabularx}\normalfont\normalsize%
\label{tab:reversion}

% -----------------------------------------------------------------------------
\section{Desmantelamiento Seguro del Entorno}%
\label{sec:desmantelamiento}

El desmantelamiento del lado del servidor se realiza con los subcomandos de Docker Compose.
Distinga entre una \textbf{detención reversible} y una \textbf{purga definitiva}.\par

\subsection[Purga de Contenedores, Redes y Volúmenes]{Detención y Purga de Contenedores, Redes Aisladas y Volúmenes de Docker}%
\label{subsec:purga-docker}

Para detener el servicio sin destruir datos —por ejemplo, durante un mantenimiento— basta
con \comando{stop}. Para una reversión completa, use \comando{down} con la opción de
volúmenes.\par

\begin{lstlisting}[language=bash]
# 1. Detencion reversible (conserva contenedores y volumenes)
$ docker compose stop

# 2. Purga: elimina contenedores y la red aislada (conserva volumenes)
$ docker compose down

# 3. Purga TOTAL: ademas elimina los volumenes nombrados (spring-cache)
$ docker compose down -v

# 4. Elimina la imagen del monolito
$ docker rmi ethoshub:2.0.0
\end{lstlisting}

\begin{PuntoNoRetorno}
El comando \comando{docker compose down -v} \textbf{elimina los volúmenes nombrados},
incluida la caché \comando{spring-cache}. Es una operación \textbf{irreversible}: la caché
de dependencias Maven se reconstruirá desde cero en el siguiente despliegue. Ejecútelo solo
si pretende una limpieza total del entorno.
\end{PuntoNoRetorno}

\par\vspace{0.2cm}
Verifique que no queden recursos residuales tras la purga:\par

\begin{lstlisting}[language=bash]
$ docker compose ps          # no debe listar contenedores
$ docker volume ls | grep ethoshub   # vacio tras 'down -v'
$ docker network ls | grep ethoshub-net  # vacio tras 'down'
\end{lstlisting}

\begin{VerificacionOK}
Si \comando{docker compose ps} no muestra contenedores y los listados de volúmenes y redes
no contienen recursos \comando{ethoshub}, el desmantelamiento del servidor fue completo. El
puerto 8080 queda libre, lo que puede confirmarse con \comando{ss -ltnp | grep 8080}.
\end{VerificacionOK}

\begin{NotaConfiguracion}
Si desplegó sin Docker (ejecución directa del JAR), la «desinstalación» consiste
simplemente en detener el proceso Java (\comando{Ctrl+C} o \comando{kill <pid>}) y, si lo
desea, borrar el directorio \comando{target/} y los archivos de log. No quedan servicios de
sistema instalados, pues el JAR es autocontenido.
\end{NotaConfiguracion}

% -----------------------------------------------------------------------------
\section{Limpieza del Almacenamiento Local de Datos del Perfil}%
\label{sec:limpieza-cliente}

EthosHub mantiene estado en el navegador del usuario: la sesión vive en
\comando{sessionStorage} y ciertas preferencias en \comando{localStorage}. Tras una
evaluación, conviene purgar este estado en los equipos compartidos.\par

\subsection[Limpieza de localStorage y sessionStorage]{Depuración de Tokens de Sesión e Historiales en \\ \texttt{localStorage} y \texttt{sessionStorage}}%
\label{subsec:limpieza-storage}

La limpieza puede hacerse desde la interfaz del navegador (ajustes de privacidad $\to$
borrar datos del sitio) o mediante la consola de desarrollo. Esta última es la más directa
y específica del dominio:\par

\begin{lstlisting}[language=bash]
// Consola de desarrollo del navegador (F12), en el dominio de EthosHub:
localStorage.clear();      // borra preferencias persistentes
sessionStorage.clear();    // borra la sesion (token JWT, estado por pestana)
// Recargue la pagina para aplicar:
location.reload();
\end{lstlisting}

\par\vspace{0.2cm}
La tabla~\ref{tab:storage} indica qué se elimina con cada almacén.\par

\vspace{0.2cm}
\noindent
\renewcommand{\arraystretch}{1.35}
\fontsize{8.7}{10.4}\selectfont\sffamily
\begin{tabularx}{\dimexpr\textwidth-2pt\relax}{|>{\ttfamily\bfseries\color{colorPrimario}}l|X|}
\hline
\rowcolor{headerTabla}
\sffamily\textbf{Almacén} & \sffamily\textbf{Contenido eliminado} \\
\hline
sessionStorage & Token de sesión (JWT) y estado de la pestaña; el cierre de sesión se propaga vía \texttt{BroadcastChannel}. \\
\hline
\rowcolor{grisTecnico}
localStorage & Preferencias persistentes (idioma, ajustes de interfaz). \\
\hline
\end{tabularx}\normalfont\normalsize%
\label{tab:storage}

\begin{AlertaSeguridad}
En equipos \textbf{compartidos} de evaluación, la limpieza del almacenamiento es
\textbf{obligatoria} al finalizar: un \textit{token} de sesión olvidado en
\comando{sessionStorage} podría permitir a otro usuario continuar la sesión. Cierre sesión
explícitamente y, a continuación, purgue el almacenamiento del navegador.
\end{AlertaSeguridad}

\begin{AvisoNota}
Si durante la evaluación aplicó las excepciones de origen inseguro del
apartado~\ref{sec:pwa-http}, recuerde \textbf{revertir también esas banderas} del
navegador: vuelva a poner en \textit{Disabled} la opción «Insecure origins treated as
secure» (Chromium) o en \comando{false} las claves \comando{media.*.insecure.enabled}
(Firefox). Así restaura la postura de seguridad original del equipo.
\end{AvisoNota}

% -----------------------------------------------------------------------------
\section{Respaldo Previo y Reversión de Versión}%
\label{sec:rollback-version}

Antes de una purga total o de actualizar a una nueva versión, conserve los elementos cuya
pérdida no es trivial. La tabla~\ref{tab:backup} indica qué respaldar y cómo.\par

\vspace{0.2cm}
\noindent
\renewcommand{\arraystretch}{1.35}
\fontsize{8.7}{10.4}\selectfont\sffamily
\begin{tabularx}{\dimexpr\textwidth-2pt\relax}{|>{\bfseries\color{colorPrimario}}l|X|}
\hline
\rowcolor{headerTabla}
\textbf{Elemento} & \textbf{Procedimiento de respaldo} \\
\hline
Base de datos & Use los \textit{backups} gestionados de Supabase (panel) o \texttt{pg\_dump} contra el \textit{pooler}. \\
\hline
\rowcolor{grisTecnico}
Archivos \texttt{uploads/} & Copie el directorio \ruta{ethoshubB/uploads/} a un almacén seguro. \\
\hline
Artefacto anterior & Conserve el JAR de la versión previa para poder revertir. \\
\hline
\rowcolor{grisTecnico}
Variables de entorno & Resguarde el archivo \texttt{.env} (cifrado) usado en el despliegue. \\
\hline
\end{tabularx}\normalfont\normalsize%
\label{tab:backup}

\par\vspace{0.2cm}
La \textbf{reversión de versión} (\textit{rollback}) consiste en volver a desplegar el
artefacto anterior. Como la persistencia reside en Supabase, el \textit{rollback} del código
es directo, pero requiere atención a las migraciones:\par

\begin{enumerate}[noitemsep]
    \item Detenga la versión actual (\comando{docker compose down} o \comando{systemctl stop ethoshub}).
    \item Despliegue el JAR (o la imagen) de la versión previa conservada.
    \item Confirme con \comando{/actuator/health} y la lista Go-Live (apartado~\ref{sec:golive}).
\end{enumerate}

\begin{AvisoAdvertencia}
Las migraciones de base de datos \textbf{no siempre son reversibles}. Si la versión nueva
aplicó una migración que la versión previa no comprende, revertir solo el código puede dejar
el esquema por delante del artefacto. Ante un \textit{rollback} que cruce migraciones,
coordine con el DBA la reversión del esquema o mantenga la compatibilidad hacia atrás.
\end{AvisoAdvertencia}

\par\vspace{0.3cm}
Con el desmantelamiento del servidor y la limpieza del cliente completados, el ciclo de
vida de la instalación —desde los requisitos del capítulo~\ref{ch:requisitos} hasta la
reversión— queda íntegramente cubierto por este manual. El operador dispone así de un
procedimiento reproducible, verificable y reversible, conforme a los principios del
estándar ISO/IEC/IEEE 15289:2024.\par
